\documentclass[11pt,twocolumn]{article}
\usepackage[margin=0.85in]{geometry}
\usepackage{amsmath,amssymb}
\usepackage{graphicx}
\usepackage{booktabs}
\usepackage{multirow}
\usepackage{hyperref}
\usepackage{xcolor}
\usepackage{enumitem}
\usepackage{algorithm}
\usepackage{algorithmic}
\usepackage{balance}

\hypersetup{colorlinks=true, linkcolor=blue!60!black, citecolor=blue!60!black, urlcolor=blue!60!black}

% Compact spacing
\setlength{\parskip}{0.3em}
\setlength{\abovedisplayskip}{6pt}
\setlength{\belowdisplayskip}{6pt}

% Citation placeholder
\newcommand{\citep}[1]{[\textcolor{red}{#1}]}

\title{\textbf{OPAL: One-Pass Approximation of Light via Physics-Embedded Flow Matching with Complex-Valued Networks}}

\author{
% Authors here
}

\date{}

\begin{document}
\maketitle

% ============================================================================
\begin{abstract}
% ============================================================================
Designing photonic integrated circuits requires accurate electromagnetic field simulations, which remain computationally expensive even for simple two-dimensional geometries. We present OPAL (One-Pass Approximation of Light), a generative neural surrogate that replaces iterative FDTD solves with a single forward pass to predict complex-valued electromagnetic field distributions for photonic devices given their geometry and operating wavelength. Our approach combines three key ideas: (1) conditional flow matching as the generative framework, learning a velocity field that transports Gaussian noise to physically valid field solutions; (2) a complex-valued U-Net architecture with native complex convolutions, magnitude-gated activations (ModReLU), and complex-domain attention that preserves phase information as a first-class quantity; and (3) a physics-constrained training curriculum that sequentially introduces a Helmholtz residual loss enforcing $\nabla^2 E_z + k_0^2 \varepsilon E_z = 0$, followed by S-parameter supervision at device ports. We introduce an interface-aware masking scheme for the Helmholtz residual that excludes dielectric boundary pixels where finite-difference stencil errors dominate, reducing the ground-truth residual floor by an order of magnitude and yielding a physically meaningful compliance metric. Trained on a dataset of 10,000 unique geometries across five device types---straight waveguides, multimode interferometers, S-bends, Y-branches, and directional couplers---simulated at three C-band wavelengths, OPAL produces field predictions in milliseconds that would otherwise require minutes of finite-difference time-domain simulation.
\end{abstract}

% ============================================================================
\section{Introduction}
\label{sec:intro}
% ============================================================================

Photonic integrated circuits (PICs) underpin modern telecommunications, sensing, quantum information, and emerging optical computing architectures \citep{Bogaerts2020, Shastri2021}. Their design relies on solving Maxwell's equations across complex dielectric geometries to predict how light propagates, couples between waveguides, and scatters at junctions. The standard simulation tool---finite-difference time-domain (FDTD) analysis---is accurate but slow, often requiring minutes to hours per geometry depending on domain size and convergence criteria. This computational cost limits design-space exploration, multi-parameter optimization, and the tight iteration cycles needed for inverse design \citep{Molesky2018, Piggott2015}.

Data-driven surrogate models offer a path to orders-of-magnitude speedup. Convolutional neural networks and U-Nets have been applied to predict electromagnetic fields directly from permittivity maps \citep{Wiecha2021, Tahersima2019}, but purely data-driven approaches ignore the governing physics, leading to solutions that may violate Maxwell's equations and generalize poorly outside the training distribution. Physics-informed neural networks (PINNs) embed PDE constraints into the loss function \citep{Raissi2019}, but typically require per-geometry optimization and do not amortize across a family of designs. Neural operators such as the Fourier Neural Operator (FNO) \citep{Li2021FNO} learn mappings between function spaces but are predominantly real-valued, discarding the phase information that is critical for coherent photonic devices where interference effects determine functionality.

A further challenge is that electromagnetic fields are inherently complex-valued. The field $E_z(x,y) = |E_z| e^{j\phi}$ carries both amplitude and phase, and the relationship between geometry and phase response is what makes devices like directional couplers, Mach-Zehnder interferometers, and multimode interferometers function. Standard real-valued neural networks that predict real and imaginary parts independently lack the inductive bias to preserve the algebraic structure of complex multiplication, which governs wave propagation, interference, and mode coupling.

In this work, we present \textbf{OPAL} (One-Pass Approximation of Light), a physics-embedded generative model for photonic field prediction. Our contributions are:

\begin{enumerate}[leftmargin=*, itemsep=2pt]
    \item \textbf{Complex-valued flow matching.} We apply conditional flow matching \citep{Lipman2023, Liu2023} with a native complex-valued U-Net to electromagnetic field prediction. The model learns a velocity field that transports Gaussian noise to valid field solutions, conditioned on device geometry and wavelength.

    \item \textbf{Helmholtz-constrained training with interface masking.} We enforce the scalar Helmholtz equation $\nabla^2 E_z + k_0^2 \varepsilon E_z = 0$ as a soft constraint during training. We introduce an interface-aware masking scheme that excludes pixels at dielectric boundaries where the finite-difference Laplacian stencil produces spurious residuals, reducing the ground-truth floor from ${\sim}0.17$ to ${\sim}0.01$--$0.03$ and providing a meaningful physics compliance metric.

    \item \textbf{Multi-loss curriculum.} We train in three phases---flow matching only, then adding physics constraints, then S-parameter supervision---with independent warmup schedules per loss. This avoids gradient conflicts that arise from simultaneous multi-objective optimization.

    \item \textbf{Multi-device generalization.} A single model handles five photonic device types with 2--5 design parameters and 2--4 ports, trained on a dataset of 30,000 FDTD simulations (10,000 geometries $\times$ 3 wavelengths) generated via Latin hypercube sampling.
\end{enumerate}

% ============================================================================
\section{Background}
\label{sec:background}
% ============================================================================

\subsection{Electromagnetic Simulation of Photonic Devices}

We consider two-dimensional TE-polarized electromagnetic fields in silicon-on-insulator (SOI) waveguides. In the frequency domain, the out-of-plane electric field component $E_z(x,y)$ satisfies the scalar Helmholtz equation:
\begin{equation}
    \nabla^2 E_z(x,y) + k_0^2 \, \varepsilon(x,y) \, E_z(x,y) = 0
    \label{eq:helmholtz}
\end{equation}
where $k_0 = 2\pi/\lambda$ is the free-space wavenumber and $\varepsilon(x,y)$ is the spatially varying relative permittivity. For SOI at C-band wavelengths ($\lambda \in [1.50, 1.60]\,\mu$m), the silicon waveguide core has an effective index $n_\text{eff} \approx 2.4$ ($\varepsilon_\text{core} \approx 5.8$) and the silica cladding has $n_\text{clad} = 1.444$ ($\varepsilon_\text{clad} \approx 2.09$).

FDTD methods solve the time-domain Maxwell's equations on a Yee grid and extract frequency-domain fields via discrete Fourier transforms \citep{Taflove2005}. While accurate, each simulation requires iterating until fields decay below a convergence threshold, typically taking minutes for the domain sizes considered here.

The key simulation outputs are: the complex field $E_z(x,y)$, the permittivity map $\varepsilon(x,y)$, and the scattering parameters (S-parameters) $S_{ij}$ that characterize port-to-port transmission and reflection.

\subsection{Conditional Flow Matching}

Flow matching \citep{Lipman2023, Liu2023} is a generative modeling framework that learns a time-dependent velocity field $u_\theta(x_t, t)$ transporting samples from a simple prior $p_0$ (typically Gaussian) to the data distribution $p_1$ along a deterministic path. Given source samples $x_0 \sim \mathcal{N}(0, I)$ and target samples $x_1 \sim p_\text{data}$, the conditional interpolation path is:
\begin{equation}
    x_t = (1 - t)\,x_0 + t\,x_1, \quad t \in [0, 1]
    \label{eq:interpolation}
\end{equation}
with corresponding target velocity $u_t = x_1 - x_0$. The training loss minimizes:
\begin{equation}
    \mathcal{L}_\text{FM} = \mathbb{E}_{t, x_0, x_1} \left[ \left\| u_\theta(x_t, t) - u_t \right\|^2 \right]
    \label{eq:fm_loss}
\end{equation}

At inference, new samples are generated by integrating the learned velocity field from $t=0$ to $t=1$ using an ODE solver, starting from Gaussian noise. Compared to diffusion models, flow matching offers simpler training dynamics, deterministic sampling, and straighter transport paths that require fewer integration steps.

% ============================================================================
\section{Method}
\label{sec:method}
% ============================================================================

\subsection{Problem Formulation}

Given a permittivity map $\varepsilon(x,y) \in \mathbb{R}^{H \times W}$, a source mask $s(x,y) \in \{0,1\}^{H \times W}$ indicating the excitation port, and a wavelength $\lambda$, we seek to predict the steady-state complex electric field $E_z(x,y) = E_r(x,y) + j\,E_i(x,y) \in \mathbb{C}^{H \times W}$.

We represent the field as a two-channel real tensor $[E_r, E_i] \in \mathbb{R}^{2 \times H \times W}$ and treat $\varepsilon$ and $s$ as conditioning inputs. The model also predicts, as an auxiliary output, the S-parameters $S_{ij} \in \mathbb{C}$ extracted from the predicted field at each port.

\subsection{Complex-Valued Physics U-Net}

\subsubsection{Architecture Overview}

Our backbone is a U-Net encoder-decoder with skip connections, operating natively in the complex domain. The encoder downsamples through a configurable number of resolution levels with channel multipliers $(1, 2, 4, 8)$ relative to a base channel count. Each level contains $N_\text{res}$ complex residual blocks, optionally followed by complex multi-head self-attention. The decoder mirrors the encoder with $N_\text{res}+1$ blocks per level. A bottleneck of two complex residual blocks with an attention layer connects encoder and decoder.

The two-channel field input $[E_r, E_i]$ is treated as a single complex channel $E_r + jE_i$. Conditioning maps (permittivity, source mask) are processed by a lightweight auxiliary encoder and concatenated with the complex features at the stem.

\subsubsection{Complex Convolutions}

Standard real-valued convolutions applied independently to real and imaginary parts cannot express the algebraic structure of complex multiplication, which is fundamental to wave physics (phase rotation, interference). We use complex convolutions \citep{Trabelsi2018} that implement:
\begin{align}
    y_r &= W_r * x_r - W_i * x_i \\
    y_i &= W_r * x_i + W_i * x_r
\end{align}
where $W_r, W_i$ are learned real-valued kernel pairs and $*$ denotes spatial convolution. This parameterization ensures that the network can represent arbitrary complex linear maps, including pure phase rotations $e^{j\phi}$ that are ubiquitous in wave propagation.

Weights are initialized using complex Glorot normal initialization with variance $1/(f_\text{in} + f_\text{out})$ for both real and imaginary kernels, where $f_\text{in/out} = C_\text{in/out} \times k^2$.

\subsubsection{ModReLU Activation}

Standard ReLU applied independently to real and imaginary parts destroys phase information. We use the ModReLU activation \citep{Arjovsky2016} which gates the magnitude while preserving the phase:
\begin{equation}
    \text{ModReLU}(z) = \text{ReLU}(|z| + b) \cdot \frac{z}{|z|}
    \label{eq:modrelu}
\end{equation}
where $b$ is a learnable bias initialized to $-0.1$. This activation zeros out signals below a learned magnitude threshold while preserving the phase angle of signals that pass through, providing a physically meaningful nonlinearity for wave fields.

\subsubsection{Complex Group Normalization}

We normalize complex features jointly rather than treating real and imaginary parts independently. Complex group normalization computes:
\begin{equation}
    \hat{z} = \frac{z - \mu_z}{\sqrt{\sigma^2_r + \sigma^2_i + \epsilon}}
\end{equation}
where $\mu_z$ is the complex mean and $\sigma^2_r, \sigma^2_i$ are the variances of real and imaginary parts within each group. The affine transform uses complex scale $(\gamma_r + j\gamma_i)$ and shift $(\beta_r + j\beta_i)$:
\begin{equation}
    y = (\gamma_r + j\gamma_i) \hat{z} + (\beta_r + j\beta_i)
\end{equation}

\subsubsection{Complex Self-Attention}

At specified resolution levels, we apply multi-head self-attention in the complex domain. Queries, keys, and values are computed via complex convolutions. The attention logits use the real part of the Hermitian inner product:
\begin{equation}
    A_{nm} = \text{softmax}\!\left(\frac{\text{Re}(Q_n \cdot \overline{K_m})}{\sqrt{d_h}}\right)
\end{equation}
where $\overline{K}$ denotes complex conjugation and $d_h$ is the head dimension. This computes similarity based on both amplitude and phase alignment. The output projection is zero-initialized for training stability.

\subsubsection{Timestep Conditioning}

The flow matching time $t \in [0,1]$ is encoded via sinusoidal positional embedding followed by an MLP. The resulting embedding is injected into each residual block via complex modulation:
\begin{align}
    h'_r &= \alpha_r \cdot h_r - \alpha_i \cdot h_i + h_r \\
    h'_i &= \alpha_r \cdot h_i + \alpha_i \cdot h_r + h_i
\end{align}
where $(\alpha_r, \alpha_i)$ are the real and imaginary modulation scales from the time embedding MLP, which is zero-initialized so that at the start of training the modulation acts as an identity.

\subsubsection{Physics Feature Embedding}

Optionally, the network computes a Helmholtz residual feature map as an additional input channel. Given the current field estimate $E_z$ and conditioning $\varepsilon$, the residual $R = \nabla^2 E_z + k_0^2 \varepsilon E_z$ is computed via finite differences and concatenated with the input features. This provides the network with an explicit signal of where its current prediction violates physics, forming a feedback loop during both training and inference.

\subsection{Training Objective}

The total loss combines four terms with a phased curriculum:
\begin{equation}
    \mathcal{L} = \mathcal{L}_\text{FM} + \lambda_R \mathcal{L}_R + \lambda_E \mathcal{L}_E + \lambda_S \mathcal{L}_S
    \label{eq:total_loss}
\end{equation}

\subsubsection{Flow Matching Loss}

The core training signal is the conditional flow matching objective (Eq.~\ref{eq:fm_loss}), applied only to the field channels. The conditioning inputs ($\varepsilon$, source mask) are concatenated with the noised field $x_t$ but are not noised themselves. Time is sampled from a 50/50 mixture of uniform $t \sim U[0,1]$ and logit-normal $t = \sigma(\mu + s\xi)$, $\xi \sim \mathcal{N}(0,1)$, which concentrates training signal near $t=0.5$ while maintaining coverage at the boundaries.

The loss is spatially weighted by a mask that:
(a) zeros out PML absorbing boundary regions,
(b) upweights device regions where fields are most complex, and
(c) ensures a minimum weight $w_\text{min}$ everywhere in the interior to prevent ignoring cladding fields.

\subsubsection{Helmholtz Residual Loss}
\label{sec:helmholtz}

The Helmholtz residual penalizes violations of Eq.~\ref{eq:helmholtz} in the model's predicted field. Given the endpoint estimate $\hat{x}_1$ (obtained by rearranging the flow matching velocity prediction), we denormalize to physical units and compute:
\begin{equation}
    R = \nabla^2 \hat{E}_z + k_0^2 \varepsilon \hat{E}_z
\end{equation}
using a 5-point finite-difference Laplacian with anisotropic grid spacing $(\Delta x, \Delta y)$:
\begin{equation}
    \nabla^2 f \approx \frac{f_{i,j+1} - 2f_{i,j} + f_{i,j-1}}{\Delta x^2} + \frac{f_{i+1,j} - 2f_{i,j} + f_{i-1,j}}{\Delta y^2}
\end{equation}

The residual is masked to exclude: (i) PML regions where the Helmholtz equation does not hold, (ii) source regions where a driving term is present, and (iii) dielectric interface pixels, as described below.

\paragraph{Interface-aware masking.}
At sharp dielectric boundaries (e.g., Si/SiO$_2$ interfaces), the normal component of $E_z$ has discontinuous derivatives. The 5-point Laplacian stencil, which assumes smooth fields, produces $O(1)$ spurious residuals at these pixels regardless of field accuracy. We identify interface pixels as those where the permittivity gradient magnitude exceeds a threshold:
\begin{equation}
    m_\text{interface}(x,y) = \mathbf{1}\!\left[|\nabla \varepsilon(x,y)| < \tau\right]
\end{equation}
where $\nabla\varepsilon$ is computed using central differences and $\tau$ is a hyperparameter (we use $\tau = 1.0$). The full residual weight mask is:
\begin{equation}
    w_R = w_\text{FM} \cdot m_\text{source-free} \cdot m_\text{interface}
\end{equation}

Without interface masking, ground-truth FDTD fields exhibit a residual ratio of ${\sim}0.17$ (i.e., 17\% of the driving term magnitude)---a floor that no model can surpass. With masking, this floor drops to ${\sim}0.01$--$0.03$, yielding a metric that meaningfully measures physics compliance.

\paragraph{Driving-term normalization.}
To make the residual loss dimensionless, grid-independent, and field-magnitude-independent, we normalize by the driving-term scale:
\begin{equation}
    \mathcal{L}_R = \frac{\langle w_R \cdot |R|^2 \rangle}{\langle w_R \cdot |k_0^2 \varepsilon E_z^\text{true}|^2 \rangle}
\end{equation}
where $\langle \cdot \rangle$ denotes spatial averaging. This ratio is approximately 0 when the Helmholtz equation is satisfied and approximately 1 when the residual is as large as the dominant PDE term.

\paragraph{Time gating.}
At early flow times ($t \approx 0$), the model's field estimate $\hat{x}_1$ is mostly noise, making the Helmholtz residual meaningless. We gate physics losses to only contribute when $t \geq t_\text{min}$ (typically $0.5$), and optionally weight each sample's residual by $t^p$ (with $p=4$) to further upweight clean, high-$t$ reconstructions.

\subsubsection{Endpoint Loss}

A direct mean squared error between the predicted endpoint $\hat{x}_1$ and the ground-truth field $x_1$, applied after global phase alignment:
\begin{equation}
    \mathcal{L}_E = \left\| \hat{x}_1 - e^{j\phi_\text{align}} x_1 \right\|^2
\end{equation}
where $\phi_\text{align}$ is the phase that maximizes the weighted inner product $\langle w \cdot \overline{\hat{E}} \cdot E_\text{true} \rangle$, resolving the global phase ambiguity inherent in frequency-domain fields.

\subsubsection{S-Parameter Loss}

S-parameters quantify port-to-port transmission and are the primary figures of merit for photonic device design. We extract S-parameters from predicted fields via modal decomposition: at each port, we solve the one-dimensional eigenmode equation along the waveguide cross-section to obtain the fundamental mode profile $\psi_m(r_\perp)$, then compute the overlap integral:
\begin{equation}
    a_p = \int E_z(r_\perp)\, \overline{\psi_m(r_\perp)}\, dr_\perp
\end{equation}
The S-parameter from input port $i$ to output port $j$ is $S_{ji} = a_j / a_i$. The loss penalizes both magnitude and phase errors:
\begin{equation}
    \mathcal{L}_S = \sum_{j} g_j \left[ \left(|\hat{S}_{ji}| - |S_{ji}^\text{true}|\right)^2 + \alpha\left(1 - \text{Re}\!\left(\frac{\hat{S}_{ji}}{|\hat{S}_{ji}|} \cdot \overline{\frac{S_{ji}^\text{true}}{|S_{ji}^\text{true}|}}\right)\right) \right]
\end{equation}
where $g_j$ is an amplitude gate that suppresses loss for near-zero transmissions (below $10^{-3}$) and excludes the input port from supervision.

\subsection{Multi-Loss Curriculum}
\label{sec:curriculum}

Introducing all losses simultaneously leads to gradient conflicts that destabilize training. We use a three-phase curriculum:

\begin{itemize}[leftmargin=*, itemsep=2pt]
    \item \textbf{Phase A} (epochs $0$ to $N_A$): Flow matching loss only. The model learns basic field structure and the velocity prediction task.
    \item \textbf{Phase B} (epochs $N_A$ to $N_B$): Add Helmholtz residual and endpoint losses with linear warmup ramps. Physics constraints guide the model toward valid solutions.
    \item \textbf{Phase C} (epochs $N_B$ onward): Add S-parameter loss. The model fine-tunes port coupling accuracy while maintaining field quality and physics compliance.
\end{itemize}

Each loss $\lambda_i(e)$ follows an independent linear ramp from 0 to its target weight over a configurable number of warmup epochs after its phase starts. This prevents gradient shock when new objectives are introduced.

\subsection{Inference}

At test time, we sample $x_0 \sim \mathcal{N}(0, I)$ and integrate the learned velocity field from $t=0$ to $t=1$ using Heun's method (second-order Runge-Kutta):
\begin{align}
    v_0 &= u_\theta(x_k, t_k) \\
    x_\text{Euler} &= x_k + \Delta t \cdot v_0 \\
    v_1 &= u_\theta(x_\text{Euler}, t_{k+1}) \\
    x_{k+1} &= x_k + \frac{\Delta t}{2}(v_0 + v_1)
\end{align}

We use a quadratic time grid $t_k = (k/N)^2$ that allocates more steps near $t=0$ where the velocity field changes most rapidly. Typical inference uses $N=100$ steps. We evaluate using the exponential moving average (EMA) of model weights with decay 0.999.

\subsection{Data Augmentation}

We exploit the vertical mirror symmetry of the elongated photonic devices in our dataset. A horizontal flip (D2 symmetry) is applied stochastically during training, effectively doubling the training data. Full D4 augmentation (rotations) is not applicable because the rectangular simulation domain and propagation direction break rotational symmetry.

% ============================================================================
\section{Dataset}
\label{sec:dataset}
% ============================================================================

\subsection{Device Types}

We generate training data for five photonic device families, summarized in Table~\ref{tab:devices}. Together, these span the essential building blocks of photonic integrated circuits: straight propagation, lateral offset routing, power splitting, and evanescent coupling.

\begin{table*}[t]
\centering
\caption{Summary of device types in the OPAL dataset.}
\label{tab:devices}
\begin{tabular}{@{}llcccl@{}}
\toprule
\textbf{Device} & \textbf{Description} & \textbf{Ports} & \textbf{Params} & \textbf{Count} & \textbf{Key Parameter Ranges} \\
\midrule
Straight waveguide & Constant-width reference & 2 & 2 & 250 & width: 0.39--0.58\,$\mu$m, length: 6--22\,$\mu$m \\
2$\times$2 MMI & Multimode interferometer & 4 & 5 & 3{,}500 & MMI width: 2.5--6.0\,$\mu$m, length: 8--20\,$\mu$m \\
S-bend & Euler-curve lateral offset & 2 & 3 & 1{,}000 & offset: 2--7\,$\mu$m, $R_\text{min}$: 3--8\,$\mu$m \\
Y-branch & 1$\times$2 power splitter & 3 & 5 & 1{,}750 & junction: 1.2--3.2\,$\mu$m, bend: 4--7\,$\mu$m \\
Dir.\ coupler & 2$\times$2 evanescent coupler & 4 & 5 & 3{,}500 & gap: 0.11--0.33\,$\mu$m, length: 5--12\,$\mu$m \\
\midrule
\multicolumn{4}{l}{\textbf{Total unique geometries}} & \textbf{10{,}000} & \\
\bottomrule
\end{tabular}
\end{table*}

All devices use silicon-on-insulator waveguides with effective core index $n_\text{eff} \approx 2.4$ (from pre-computed eigenmode tables for 220\,nm slab thickness) and SiO$_2$ cladding with $n_\text{clad} = 1.444$. Waveguide widths range from 0.39 to 0.58\,$\mu$m across all devices, supporting single-mode operation in the C-band.

\subsection{FDTD Simulation}

Each device geometry is simulated at three wavelengths: $\lambda \in \{1.50, 1.55, 1.60\}\,\mu$m, producing 30,000 total simulations. We use the Meep FDTD solver \citep{Oskooi2010} with the following configuration:

\begin{itemize}[leftmargin=*, itemsep=2pt]
    \item \textbf{Resolution}: 18 pixels/$\mu$m ($\Delta x = 55.6$\,nm).
    \item \textbf{PML}: 0.67\,$\mu$m absorbing boundary on all sides.
    \item \textbf{Source}: Eigenmode source exciting the fundamental TE mode at a randomly selected input port.
    \item \textbf{Convergence}: Field decay threshold of $10^{-5}$.
    \item \textbf{Interior crop}: $512 \times 192$ pixels ($28.4 \times 10.7$\,$\mu$m), excluding PML regions.
\end{itemize}

Per simulation, we store: complex $E_z$ field, permittivity $\varepsilon$, source mask, per-port binary masks, and reference S-parameters extracted by Meep's eigenmode decomposition.

\subsection{Sampling and Splitting}

Device parameters are sampled via Latin hypercube sampling (LHS) within the ranges in Table~\ref{tab:devices}, with values quantized to a half-pixel grid ($1/(2 \times \text{resolution})$) to ensure each parameter change produces a measurably different subpixel-averaged permittivity. The allocation across device types is proportional to parameter-space dimensionality, giving more samples to the 5-parameter families (MMI, Y-branch, directional coupler).

Data is split 80/10/10 into train/validation/test by geometry ID, ensuring that all wavelengths for a given geometry belong to the same split (no wavelength leakage). Data is stored in sharded compressed NPZ files with a JSON index for efficient loading.

\subsection{Phase Anchoring}

Frequency-domain fields have an arbitrary global phase: if $E_z$ is a solution, so is $E_z e^{j\phi}$ for any constant $\phi$. To provide a consistent training target, we anchor the global phase using the field within the source region. We compute a weighted average of unit phasors in the source mask weighted by local field intensity:
\begin{equation}
    \phi_\text{ref} = \arg\!\left(\sum_{(x,y) \in \text{src}} |E_z(x,y)|^2 \cdot \frac{E_z(x,y)}{|E_z(x,y)|}\right)
\end{equation}
and rotate the entire field by $e^{-j\phi_\text{ref}}$, ensuring the source region has approximately zero phase. If the source mask is unavailable or too small, we fall back to a fixed region-of-interest near the input waveguide.

% ============================================================================
\section{Experiments}
\label{sec:experiments}
% ============================================================================

\subsection{Training Configuration}

We train on 2 NVIDIA A100 GPUs using distributed data parallel (DDP). The model uses a base channel count of 32, channel multipliers $(1, 2, 4, 8)$, 3 residual blocks per level, and 8-head attention at the coarsest resolution ($8\times$ downsampled). We use AdamW \citep{Loshchilov2019} with learning rate $10^{-4}$, weight decay $5 \times 10^{-5}$, and a schedule combining 10-epoch linear warmup with cosine decay to $5 \times 10^{-6}$. Mixed precision training (FP16 forward pass, FP32 losses) accelerates training while maintaining numerical stability in physics computations. Batch size is 60 across both GPUs. EMA with decay 0.999 is applied at every step.

The curriculum is configured as: Phase A for 10 epochs (FM only), Phase B from epoch 10 (add residual $\lambda_R = 3.0$ with 10-epoch warmup, endpoint $\lambda_E = 0.5$ with 20-epoch warmup), and S-parameter loss starting in Phase B ($\lambda_S = 0.3$ with 50-epoch warmup) using modal decomposition. Physics losses are time-gated to $t \geq 0.5$ with residual $t$-power weighting $p=4$.

\subsection{Baselines and Ablations}

We evaluate the following configurations:

\begin{itemize}[leftmargin=*, itemsep=2pt]
    \item \textbf{Real-valued U-Net}: Same architecture but with standard real convolutions, applied independently to $E_r$ and $E_i$. Tests the importance of complex-valued inductive bias.
    \item \textbf{FM only}: Flow matching loss without physics or S-parameter constraints. Tests the value of physics-embedded training.
    \item \textbf{No curriculum}: All losses active from epoch 0. Tests whether phased introduction matters.
    \item \textbf{No interface masking}: Helmholtz residual computed without excluding interface pixels. Tests the impact of our masking scheme on training dynamics and final metrics.
\end{itemize}

\subsection{Evaluation Metrics}

\paragraph{Field quality.}
We report peak signal-to-noise ratio (PSNR), structural similarity index (SSIM), and relative $L_2$ error for both amplitude $|E_z|$ and phase $\angle E_z$ on the held-out test set. Phase metrics are computed after global phase alignment and weighted by field amplitude to suppress noise in low-signal regions.

\paragraph{Physics compliance.}
The Helmholtz residual ratio (Section~\ref{sec:helmholtz}) with interface masking, reported relative to the ground-truth floor. A ratio near 1.0$\times$ GT indicates the model satisfies the Helmholtz equation as well as the FDTD reference.

\paragraph{S-parameter accuracy.}
Mean absolute error in $|S_{ij}|$ (linear and dB scale) and phase error in degrees, broken down by device type and port.

\paragraph{Computational cost.}
Inference time per sample (single and batched) compared to FDTD wall-clock time. Model parameter count and training GPU-hours.

\subsection{Results}

\textcolor{red}{[Results to be filled in after training runs complete. Include:]}

\begin{itemize}[leftmargin=*, itemsep=2pt]
    \item \textcolor{red}{Table: PSNR, SSIM, L2 error per device type for complex vs.\ real UNet.}
    \item \textcolor{red}{Table: S-parameter magnitude and phase error per device type.}
    \item \textcolor{red}{Figure: Side-by-side GT vs.\ predicted fields (amplitude, phase, error map) for representative devices.}
    \item \textcolor{red}{Figure: Helmholtz residual maps with and without interface masking.}
    \item \textcolor{red}{Figure: Training loss curves showing curriculum phase transitions.}
    \item \textcolor{red}{Table: Ablation study results (complex vs.\ real, with/without physics, with/without curriculum).}
    \item \textcolor{red}{Table: Inference speed vs.\ FDTD.}
\end{itemize}

% ============================================================================
\section{Discussion}
\label{sec:discussion}
% ============================================================================

\paragraph{Why complex-valued networks matter.}
Photonic devices function through interference, where the relative phase between optical paths determines constructive or destructive combination. A directional coupler's splitting ratio, for instance, depends on $\sin^2(\kappa L)$ where $\kappa$ is the coupling coefficient---a quantity determined entirely by the phase relationship between the even and odd supermodes. Real-valued networks that predict $E_r$ and $E_i$ independently lack the inductive bias to preserve the algebraic structure of complex multiplication ($e^{j\phi}$ rotations, interference sums), forcing the network to learn these relationships from data alone. Complex convolutions encode this structure directly, reducing the learning burden and improving phase accuracy.

\paragraph{Interface masking as a general principle.}
The observation that finite-difference PDE residuals produce spurious errors at material discontinuities is not specific to photonics. Any physics-informed approach that evaluates a PDE residual on a grid with discontinuous coefficients will face the same issue. Our interface masking scheme---excluding pixels where $|\nabla\varepsilon|$ exceeds a threshold---is applicable to any domain with sharp material boundaries (acoustic, thermal, elastodynamic), and we believe it should be standard practice when reporting PDE residual metrics on heterogeneous media.

\paragraph{Curriculum training for multi-objective physics.}
The gradient dynamics of simultaneously optimizing flow matching velocity, Helmholtz residual, and S-parameter objectives are complex. The FM loss operates on the velocity field at arbitrary times $t$, the residual loss operates on the endpoint field estimate, and the S-parameter loss involves a non-differentiable mode-solving step followed by a ratio. These objectives can produce conflicting gradient directions, particularly early in training when the model's field estimates are poor. The phased curriculum resolves this by ensuring the model first learns a reasonable velocity field before being asked to satisfy physics constraints, and learns reasonable fields before being asked to produce accurate S-parameters.

\paragraph{Limitations.}
Our approach has several limitations. First, we operate in 2D (TE polarization only); extension to full 3D vectorial fields would require substantially more data and compute. Second, the simulation resolution of 18 pixels/$\mu$m limits the minimum resolvable feature to approximately 110\,nm (2 pixels), which constrains the range of gap sizes in directional couplers. Third, each prediction corresponds to a single frequency; broadband prediction would require either multiple forward passes or a model architecture that conditions on a continuous wavelength range. Fourth, our device types are parameterized geometries; generalization to arbitrary freeform topologies remains an open challenge. Finally, the model predicts fields for a given input port excitation; obtaining the full S-matrix requires one forward pass per input port.

% ============================================================================
\section{Related Work}
\label{sec:related}
% ============================================================================

\paragraph{Neural surrogates for electromagnetic simulation.}
Data-driven models for photonic device simulation have been explored using fully-connected networks \citep{Tahersima2019}, convolutional networks \citep{Wiecha2021}, and U-Nets \citep{Chen2020Surrogate}. These approaches typically predict field amplitudes or S-parameters directly from geometry parameters, treating the problem as a regression task. Our work differs in using a generative framework (flow matching) that naturally handles the multi-modal nature of field solutions and in operating on full spatial field maps rather than scalar outputs.

\paragraph{Physics-informed approaches.}
Physics-informed neural networks \citep{Raissi2019} embed PDE residuals into the loss function, and have been applied to electromagnetic problems including waveguide mode solving \citep{Chen2020PINN} and scattering \citep{Fang2021}. However, PINNs optimize per-instance and do not amortize across geometries. Our approach shares the PDE residual loss concept but applies it to a pre-trained generative model, achieving amortized inference across thousands of geometries.

\paragraph{Neural operators.}
The Fourier Neural Operator \citep{Li2021FNO} and DeepONet \citep{Lu2021} learn mappings between function spaces and have shown promise for PDE problems. FNO has been applied to electromagnetic problems \citep{Lim2023} but typically in the real-valued setting. Our complex-valued architecture provides a complementary inductive bias specifically suited to wave physics.

\paragraph{Generative models for scientific computing.}
Diffusion models and flow matching have been applied to molecular dynamics \citep{Hoogeboom2022}, weather prediction \citep{Price2024}, and fluid dynamics \citep{Lienen2024}. Our work extends this paradigm to electromagnetic field prediction with domain-specific physics constraints.

\paragraph{Complex-valued neural networks.}
Complex-valued networks have been explored for signal processing \citep{Hirose2012}, MRI reconstruction \citep{Cole2021}, and speech enhancement \citep{Choi2019}. The ModReLU activation was introduced by \citep{Arjovsky2016} for unitary RNNs. Our contribution is combining complex-valued architectures with flow matching and PDE constraints for photonic simulation.

% ============================================================================
\section{Conclusion}
\label{sec:conclusion}
% ============================================================================

We have presented OPAL, a physics-embedded flow matching model for predicting electromagnetic fields in photonic integrated circuit devices. By combining conditional flow matching with a complex-valued U-Net architecture, Helmholtz residual constraints with interface-aware masking, and a multi-loss training curriculum, our approach produces field predictions that respect Maxwell's equations while accurately capturing the amplitude and phase distributions that determine device functionality.

Key to making physics-constrained training effective is our observation that finite-difference PDE residuals at dielectric interfaces are dominated by discretization artifacts rather than actual physics violations. Interface masking reduces the ground-truth residual floor by an order of magnitude, enabling the Helmholtz loss to provide meaningful gradient signal during training. We believe this insight generalizes to any physics-informed approach operating on domains with material discontinuities.

\paragraph{Future work.}
Several directions remain. Extension to three-dimensional simulations would capture out-of-plane confinement effects critical for real devices. Integration with inverse design loops---using the differentiable model as a fast forward solver within gradient-based topology optimization---is a natural application. Broadband models that predict fields across a continuous wavelength range from a single forward pass would further increase utility. Finally, scaling to larger model sizes and datasets, following the trends observed in other domains, may unlock generalization to freeform geometries beyond parameterized device families.

% ============================================================================
% References
% ============================================================================
\section*{References}
\small
\textcolor{red}{[Citations to be added. Key references to include:]}

\begin{enumerate}[leftmargin=*, itemsep=1pt, label={[\arabic*]}]
    \item \textcolor{red}{Lipman et al., ``Flow Matching for Generative Modeling,'' ICLR 2023.}
    \item \textcolor{red}{Liu et al., ``Flow Straight and Fast,'' ICLR 2023.}
    \item \textcolor{red}{Raissi et al., ``Physics-informed neural networks,'' JCP 2019.}
    \item \textcolor{red}{Li et al., ``Fourier Neural Operator,'' ICLR 2021.}
    \item \textcolor{red}{Oskooi et al., ``Meep: A flexible free-software package for electromagnetic simulations by the FDTD method,'' CPC 2010.}
    \item \textcolor{red}{Trabelsi et al., ``Deep Complex Networks,'' ICLR 2018.}
    \item \textcolor{red}{Arjovsky et al., ``Unitary Evolution Recurrent Neural Networks,'' ICML 2016.}
    \item \textcolor{red}{Taflove \& Hagness, ``Computational Electrodynamics: The FDTD Method,'' 2005.}
    \item \textcolor{red}{Bogaerts et al., ``Programmable photonic circuits,'' Nature 2020.}
    \item \textcolor{red}{Molesky et al., ``Inverse design in nanophotonics,'' Nature Photonics 2018.}
    \item \textcolor{red}{Piggott et al., ``Inverse design and demonstration of a compact and broadband on-chip wavelength demultiplexer,'' Nature Photonics 2015.}
    \item \textcolor{red}{Loshchilov \& Hutter, ``Decoupled Weight Decay Regularization,'' ICLR 2019.}
    \item \textcolor{red}{Wiecha et al., ``Deep learning in nano-photonics,'' 2021.}
    \item \textcolor{red}{Tahersima et al., ``Deep neural network inverse design of integrated photonic power splitters,'' Scientific Reports 2019.}
    \item \textcolor{red}{Lu et al., ``Learning nonlinear operators via DeepONet,'' Nature Machine Intelligence 2021.}
    \item \textcolor{red}{+ Additional photonics ML, diffusion/FM for science, complex NN references.}
\end{enumerate}

% ============================================================================
\appendix
\section{Architecture Details}
% ============================================================================

\begin{table}[h]
\centering
\caption{Model architecture hyperparameters.}
\begin{tabular}{@{}ll@{}}
\toprule
\textbf{Parameter} & \textbf{Value} \\
\midrule
Base channels & 32 \\
Channel multipliers & (1, 2, 4, 8) \\
Residual blocks per level & 3 \\
Attention heads & 8 \\
Attention resolutions & 8$\times$ downsample \\
Dropout & 0.0 \\
Activation & ModReLU \\
Normalization & Complex GroupNorm (32 groups) \\
Time embedding dim & $4 \times$ base channels \\
\bottomrule
\end{tabular}
\end{table}

\section{Training Hyperparameters}

\begin{table}[h]
\centering
\caption{Training configuration.}
\begin{tabular}{@{}ll@{}}
\toprule
\textbf{Parameter} & \textbf{Value} \\
\midrule
Optimizer & AdamW (fused) \\
Learning rate & $10^{-4}$ \\
Min learning rate & $5 \times 10^{-6}$ \\
Weight decay & $5 \times 10^{-5}$ \\
LR schedule & Linear warmup + cosine decay \\
Warmup epochs & 10 \\
Batch size & 60 (global) \\
Precision & FP16 forward / FP32 loss \\
EMA decay & 0.999 \\
\midrule
Phase A epochs & 10 \\
Phase B epochs & 50 \\
$\lambda_\text{residual}$ & 3.0 (10-epoch warmup) \\
$\lambda_\text{endpoint}$ & 0.5 (20-epoch warmup) \\
$\lambda_\text{sparam}$ & 0.3 (50-epoch warmup) \\
$t_\text{physics\_min}$ & 0.5 \\
Residual $t$-power & 4.0 \\
Interface threshold $\tau$ & 1.0 \\
\midrule
Time sampling & 50/50 uniform + logit-normal \\
ODE solver & Heun (RK2) \\
ODE steps (inference) & 100 \\
Time grid & Linear \\
\bottomrule
\end{tabular}
\end{table}

\section{Device Parameter Ranges}

\begin{table}[h]
\centering
\caption{Full parameter ranges per device type. All lengths in $\mu$m. Values quantized to $1/(2 \times 18) \approx 0.028\,\mu$m grid.}
\small
\begin{tabular}{@{}lll@{}}
\toprule
\textbf{Device} & \textbf{Parameter} & \textbf{Range} \\
\midrule
\multirow{2}{*}{Straight} & wg\_width & 0.389--0.583 \\
& dev\_length & 6.0--22.0 \\
\midrule
\multirow{5}{*}{MMI} & wg\_width & 0.389--0.583 \\
& mmi\_width & 2.5--6.0 \\
& mmi\_length & 8.0--20.0 \\
& taper\_width & 0.583--1.5 \\
& taper\_length & 1.0--3.0 \\
\midrule
\multirow{3}{*}{S-bend} & wg\_width & 0.389--0.583 \\
& lateral\_offset & 2.0--7.0 \\
& R\_min & 3.0--8.0 \\
\midrule
\multirow{5}{*}{Y-branch} & wg\_width & 0.389--0.583 \\
& l\_junction & 1.167--3.167 \\
& l\_bend & 4.0--7.0 \\
& h\_bend & 0.556--3.0 \\
& l\_out & 1.0--6.0 \\
\midrule
\multirow{5}{*}{Dir.\ coupler} & wg\_width & 0.389--0.583 \\
& gap & 0.111--0.333 \\
& wg\_length & 5.0--12.0 \\
& bend\_length & 4.0--6.0 \\
& lead\_extra\_gap & 0.806--2.5 \\
\bottomrule
\end{tabular}
\end{table}

\end{document}
